\documentclass[10pt,a4paper]{amsart}
\usepackage{amsmath,amssymb,amsthm,graphicx}
% Retained preamble for the Stacks Project chapter {CH} (shared project preamble).
% The chapter file pinned-original-source.tex only carries \input{{preamble}} on its first line,
% so no preamble text lives inside the chapter itself.  The Stacks Project uses one single shared
% preamble for every chapter: the snapshot stages it three times (sources/stacks-chow/preamble.tex,
% sources/stacks-etale-cohomology/preamble.tex and focused-seeds/AGP/sources/stacks-snapshot-20261011/
% preamble.tex) and those copies are identical in size.  The block below is copied mechanically,
% line for line, from sources/stacks-chow/preamble.tex lines 14-18 and 154-217: the Xy-pic setup and
% the native theorem declarations together with the mathematical macros.  Nothing was added, removed
% or redefined; the wrapper class and the amsmath/amssymb/amsthm/graphicx base packages are supplied
% by the assembler.
% For commutative diagrams we use Xy-pic
% Chapter: stacks-spaces-more-morphisms.tex -- More on Morphisms of Spaces.
\usepackage[all]{xy}

% We use 2cell for 2-commutative diagrams.
\xyoption{2cell}

% Theorem environments.
%
\theoremstyle{plain}
\newtheorem{theorem}[subsection]{Theorem}
\newtheorem{proposition}[subsection]{Proposition}
\newtheorem{lemma}[subsection]{Lemma}

\theoremstyle{definition}
\newtheorem{definition}[subsection]{Definition}
\newtheorem{example}[subsection]{Example}
\newtheorem{exercise}[subsection]{Exercise}
\newtheorem{situation}[subsection]{Situation}

\theoremstyle{remark}
\newtheorem{remark}[subsection]{Remark}
\newtheorem{remarks}[subsection]{Remarks}

\numberwithin{equation}{subsection}

% Macros
%
\def\lim{\mathop{\mathrm{lim}}\nolimits}
\def\colim{\mathop{\mathrm{colim}}\nolimits}
\def\Spec{\mathop{\mathrm{Spec}}}
\def\Hom{\mathop{\mathrm{Hom}}\nolimits}
\def\Ext{\mathop{\mathrm{Ext}}\nolimits}
\def\SheafHom{\mathop{\mathcal{H}\!\mathit{om}}\nolimits}
\def\SheafExt{\mathop{\mathcal{E}\!\mathit{xt}}\nolimits}
\def\Sch{\mathit{Sch}}
\def\Mor{\mathop{\mathrm{Mor}}\nolimits}
\def\Ob{\mathop{\mathrm{Ob}}\nolimits}
\def\Sh{\mathop{\mathit{Sh}}\nolimits}
\def\NL{\mathop{N\!L}\nolimits}
\def\CH{\mathop{\mathrm{CH}}\nolimits}
\def\proetale{{pro\text{-}\acute{e}tale}}
\def\etale{{\acute{e}tale}}
\def\QCoh{\mathit{QCoh}}
\def\Ker{\mathop{\mathrm{Ker}}}
\def\Im{\mathop{\mathrm{Im}}}
\def\Coker{\mathop{\mathrm{Coker}}}
\def\Coim{\mathop{\mathrm{Coim}}}

% Boxtimes
%
\DeclareMathSymbol{\boxtimes}{\mathbin}{AMSa}{"02}

%
% Macros for moduli stacks/spaces
%
\def\QCohstack{\mathcal{QC}\!\mathit{oh}}
\def\Cohstack{\mathcal{C}\!\mathit{oh}}
\def\Spacesstack{\mathcal{S}\!\mathit{paces}}
\def\Quotfunctor{\mathrm{Quot}}
\def\Hilbfunctor{\mathrm{Hilb}}
\def\Curvesstack{\mathcal{C}\!\mathit{urves}}
\def\Polarizedstack{\mathcal{P}\!\mathit{olarized}}
\def\Complexesstack{\mathcal{C}\!\mathit{omplexes}}
% \Pic is the operator that assigns to X its picard group, usage \Pic(X)
% \Picardstack_{X/B} denotes the Picard stack of X over B
% \Picardfunctor_{X/B} denotes the Picard functor of X over B
\def\Pic{\mathop{\mathrm{Pic}}\nolimits}
\def\Picardstack{\mathcal{P}\!\mathit{ic}}
\def\Picardfunctor{\mathrm{Pic}}
\def\Deformationcategory{\mathcal{D}\!\mathit{ef}}
\usepackage{cleveref}
\crefname{originallemma}{original lemma}{original lemmas}
\Crefname{originallemma}{Original lemma}{Original lemmas}
\makeatletter
\newlabel{decent-spaces-lemma-base-change-relative-conditions}{{03L0}{source}{}{sourcepointer-1}{}}
\newlabel{decent-spaces-lemma-base-change-relative-conditions@cref}{{[originallemma][0][]03L0}{[1][0][]source}{}{}{}}
\makeatother
\crefname{originallemma}{original lemma}{original lemmas}
\Crefname{originallemma}{Original lemma}{Original lemmas}
\makeatletter
\newlabel{decent-spaces-lemma-composition-relative-conditions}{{03L1}{source}{}{sourcepointer-2}{}}
\newlabel{decent-spaces-lemma-composition-relative-conditions@cref}{{[originallemma][0][]03L1}{[1][0][]source}{}{}{}}
\makeatother
\crefname{originallemma}{original lemma}{original lemmas}
\Crefname{originallemma}{Original lemma}{Original lemmas}
\makeatletter
\newlabel{decent-spaces-lemma-properties-trivial-implications}{{03M5}{source}{}{sourcepointer-3}{}}
\newlabel{decent-spaces-lemma-properties-trivial-implications@cref}{{[originallemma][0][]03M5}{[1][0][]source}{}{}{}}
\makeatother
\crefname{originallemma}{original lemma}{original lemmas}
\Crefname{originallemma}{Original lemma}{Original lemmas}
\makeatletter
\newlabel{decent-spaces-lemma-when-field}{{047Z}{source}{}{sourcepointer-4}{}}
\newlabel{decent-spaces-lemma-when-field@cref}{{[originallemma][0][]047Z}{[1][0][]source}{}{}{}}
\makeatother
\crefname{originallemma}{original lemma}{original lemmas}
\Crefname{originallemma}{Original lemma}{Original lemmas}
\makeatletter
\newlabel{definition-radicial}{{0481}{source}{}{sourcepointer-5}{}}
\newlabel{definition-radicial@cref}{{[originallemma][0][]0481}{[1][0][]source}{}{}{}}
\makeatother
\crefname{originallemma}{original lemma}{original lemmas}
\Crefname{originallemma}{Original lemma}{Original lemmas}
\makeatletter
\newlabel{morphisms-lemma-universally-injective}{{01S4}{source}{}{sourcepointer-6}{}}
\newlabel{morphisms-lemma-universally-injective@cref}{{[originallemma][0][]01S4}{[1][0][]source}{}{}{}}
\makeatother
\crefname{originallemma}{original lemma}{original lemmas}
\Crefname{originallemma}{Original lemma}{Original lemmas}
\makeatletter
\newlabel{schemes-lemma-immersion-when-closed}{{01IQ}{source}{}{sourcepointer-7}{}}
\newlabel{schemes-lemma-immersion-when-closed@cref}{{[originallemma][0][]01IQ}{[1][0][]source}{}{}{}}
\makeatother
\crefname{originallemma}{original lemma}{original lemmas}
\Crefname{originallemma}{Original lemma}{Original lemmas}
\makeatletter
\newlabel{spaces-morphisms-lemma-base-change-separated}{{03KL}{source}{}{sourcepointer-8}{}}
\newlabel{spaces-morphisms-lemma-base-change-separated@cref}{{[originallemma][0][]03KL}{[1][0][]source}{}{}{}}
\makeatother
\crefname{originallemma}{original lemma}{original lemmas}
\Crefname{originallemma}{Original lemma}{Original lemmas}
\makeatletter
\newlabel{spaces-morphisms-lemma-composition-separated}{{03KQ}{source}{}{sourcepointer-9}{}}
\newlabel{spaces-morphisms-lemma-composition-separated@cref}{{[originallemma][0][]03KQ}{[1][0][]source}{}{}{}}
\makeatother
\crefname{originallemma}{original lemma}{original lemmas}
\Crefname{originallemma}{Original lemma}{Original lemmas}
\makeatletter
\newlabel{spaces-morphisms-lemma-immersions-monomorphisms}{{042R}{source}{}{sourcepointer-10}{}}
\newlabel{spaces-morphisms-lemma-immersions-monomorphisms@cref}{{[originallemma][0][]042R}{[1][0][]source}{}{}{}}
\makeatother
\crefname{originallemma}{original lemma}{original lemmas}
\Crefname{originallemma}{Original lemma}{Original lemmas}
\makeatletter
\newlabel{spaces-properties-lemma-quasi-compact-affine-cover}{{03H6}{source}{}{sourcepointer-11}{}}
\newlabel{spaces-properties-lemma-quasi-compact-affine-cover@cref}{{[originallemma][0][]03H6}{[1][0][]source}{}{}{}}
\makeatother
\crefname{originallemma}{original lemma}{original lemmas}
\Crefname{originallemma}{Original lemma}{Original lemmas}
\makeatletter
\newlabel{spaces-properties-lemma-quasi-compact-space}{{03E4}{source}{}{sourcepointer-12}{}}
\newlabel{spaces-properties-lemma-quasi-compact-space@cref}{{[originallemma][0][]03E4}{[1][0][]source}{}{}{}}
\makeatother
\title{A universally injective morphism of algebraic spaces is radicial when it is decent, quasi-separated or locally separated}
\author{The Stacks Project Authors}
\date{stacks-project:spaces-more-morphisms.tex:local-snapshot-20261011; GNU FDL 1.2 or later}
\begin{document}
\maketitle
\noindent\textit{Editorial note.} All delivered mathematical text is the author's own, copied from the pinned Stacks Project chapter file spaces-more-morphisms.tex; no statement or proof step was rewritten, reordered or removed.\par\medskip
\noindent\textit{Editorial note.} Preamble provenance, as required for this source: the first line of the pinned chapter file is \textbackslash{}input\{preamble\}, and the native preamble was not staged with the source, so no preamble text lives inside the chapter. The retained preamble used here is the project's single Stacks-wide preamble (the same 217-line / 8116-byte file that the already published chow and etale-cohomology cards use), copied mechanically line for line from sources/stacks-chow/preamble.tex lines 14-18 (Xy-pic setup) and 154-217 (theorem declarations and mathematical macros). Not one character was changed, no package was added, no macro was defined, and the excerpt serves only as preamble, never as a span coordinate.\par\medskip
\noindent\textit{Editorial note.} The same-chapter pointers are located at their \textbackslash{}label lines inside the pinned chapter file; the cross-chapter pointers are located in the project's own tag list staged in the snapshot (sources/stacks-chow/tags.txt), at the verbatim lines carrying their canonical tags.\par\medskip
\noindent\textit{Allowed standard prerequisites (editorial).}\begin{itemize}
\item Radiciality of a morphism of algebraic spaces, tested after base change to fields. Same chapter, Definition labelled definition-radicial; invoked in the proof at lines 116-119. Fixes what has to be shown for the reduced fibre (Spec(K) x\_Y X)\_red -> Spec(K).
\item Radicial and universally injective agree for morphisms of schemes. Morphisms, Lemma labelled morphisms-lemma-universally-injective; invoked at line 132. Converts the reduced-fibre conclusion back into radiciality once the structure morphism is a morphism of schemes.
\item Stability under base change and composition, and closed immersions, for the properties decent, quasi-separated, locally separated. Decent Spaces lemmas at proof lines 136-139 and Morphisms of Spaces lemmas at proof lines 157-162. Makes the uniform first paragraph applicable to each of the three properties listed in the statement.
\item A nonempty decent reduced algebraic space universally injective over a field is the spectrum of a field. Decent Spaces, Lemma labelled decent-spaces-lemma-when-field; invoked at line 145. Supplies the field-spectrum conclusion that case (1) needs.
\end{itemize}\medskip
\ifdefined\hypertarget\hypertarget{sourcepointer-1}{}\fi
\noindent\textit{Source reference.} Original lemma 03L0: sources/stacks-chow/tags.txt, source lines 4696--4696. Navigation pointer for \textbackslash{}ref\{decent-spaces-lemma-base-change-relative-conditions\} used inside the counted span (line 137); the argument that consumes it is retained in the span.\par\medskip
\ifdefined\hypertarget\hypertarget{sourcepointer-2}{}\fi
\noindent\textit{Source reference.} Original lemma 03L1: sources/stacks-chow/tags.txt, source lines 4697--4697. Navigation pointer for \textbackslash{}ref\{decent-spaces-lemma-composition-relative-conditions\} used inside the counted span (line 138); the argument that consumes it is retained in the span.\par\medskip
\ifdefined\hypertarget\hypertarget{sourcepointer-3}{}\fi
\noindent\textit{Source reference.} Original lemma 03M5: sources/stacks-chow/tags.txt, source lines 4737--4737. Navigation pointer for \textbackslash{}ref\{decent-spaces-lemma-properties-trivial-implications\} used inside the counted span (line 139); the argument that consumes it is retained in the span.\par\medskip
\ifdefined\hypertarget\hypertarget{sourcepointer-4}{}\fi
\noindent\textit{Source reference.} Original lemma 047Z: sources/stacks-chow/tags.txt, source lines 5526--5526. Navigation pointer for \textbackslash{}ref\{decent-spaces-lemma-when-field\} used inside the counted span (line 145); the argument that consumes it is retained in the span.\par\medskip
\ifdefined\hypertarget\hypertarget{sourcepointer-5}{}\fi
\noindent\textit{Source reference.} Original lemma 0481: sources/stacks-spaces-more-morphisms/pinned-original-source.tex, source lines 59--59. Navigation pointer for \textbackslash{}ref\{definition-radicial\} used inside the counted span (line 117); the argument that consumes it is retained in the span.\par\medskip
\ifdefined\hypertarget\hypertarget{sourcepointer-6}{}\fi
\noindent\textit{Source reference.} Original lemma 01S4: sources/stacks-chow/tags.txt, source lines 2325--2325. Navigation pointer for \textbackslash{}ref\{morphisms-lemma-universally-injective\} used inside the counted span (line 132); the argument that consumes it is retained in the span.\par\medskip
\ifdefined\hypertarget\hypertarget{sourcepointer-7}{}\fi
\noindent\textit{Source reference.} Original lemma 01IQ: sources/stacks-chow/tags.txt, source lines 1987--1987. Navigation pointer for \textbackslash{}ref\{schemes-lemma-immersion-when-closed\} used inside the counted span (line 181); the argument that consumes it is retained in the span.\par\medskip
\ifdefined\hypertarget\hypertarget{sourcepointer-8}{}\fi
\noindent\textit{Source reference.} Original lemma 03KL: sources/stacks-chow/tags.txt, source lines 4681--4681. Navigation pointer for \textbackslash{}ref\{spaces-morphisms-lemma-base-change-separated\} used inside the counted span (line 160); the argument that consumes it is retained in the span.\par\medskip
\ifdefined\hypertarget\hypertarget{sourcepointer-9}{}\fi
\noindent\textit{Source reference.} Original lemma 03KQ: sources/stacks-chow/tags.txt, source lines 4686--4686. Navigation pointer for \textbackslash{}ref\{spaces-morphisms-lemma-composition-separated\} used inside the counted span (line 161); the argument that consumes it is retained in the span.\par\medskip
\ifdefined\hypertarget\hypertarget{sourcepointer-10}{}\fi
\noindent\textit{Source reference.} Original lemma 042R: sources/stacks-chow/tags.txt, source lines 5338--5338. Navigation pointer for \textbackslash{}ref\{spaces-morphisms-lemma-immersions-monomorphisms\} used inside the counted span (line 162); the argument that consumes it is retained in the span.\par\medskip
\ifdefined\hypertarget\hypertarget{sourcepointer-11}{}\fi
\noindent\textit{Source reference.} Original lemma 03H6: sources/stacks-chow/tags.txt, source lines 4555--4555. Navigation pointer for \textbackslash{}ref\{spaces-properties-lemma-quasi-compact-affine-cover\} used inside the counted span (line 170); the argument that consumes it is retained in the span.\par\medskip
\ifdefined\hypertarget\hypertarget{sourcepointer-12}{}\fi
\noindent\textit{Source reference.} Original lemma 03E4: sources/stacks-chow/tags.txt, source lines 4443--4443. Navigation pointer for \textbackslash{}ref\{spaces-properties-lemma-quasi-compact-space\} used inside the counted span (line 167); the argument that consumes it is retained in the span.\par\medskip
\begin{lemma}
\label{lemma-when-universally-injective-radicial}
Let $S$ be a scheme. Let $f : X \to Y$ be a universally injective
morphism of algebraic spaces over $S$.
\begin{enumerate}
\item If $f$ is decent then $f$ is radicial.
\item If $f$ is quasi-separated then $f$ is radicial.
\item If $f$ is locally separated then $f$ is radicial.
\end{enumerate}
\end{lemma}

\begin{proof}
Let $\mathcal{P}$ be a property of morphisms of algebraic spaces
which is stable under base change and composition and holds for
closed immersions. Assume $f : X \to Y$ has $\mathcal{P}$ and
is universally injective. Then, in the situation of
Definition \ref{definition-radicial}
the morphism $(\Spec(K) \times_Y X)_{red} \to \Spec(K)$
is universally injective and has $\mathcal{P}$. This reduces the
problem of proving
$$
\mathcal{P} + \text{universally injective}
\Rightarrow
\text{radicial}
$$
to the problem of proving that any nonempty reduced algebraic space $X$
over field whose structure morphism $X \to \Spec(K)$ is universally
injective and $\mathcal{P}$ is representable by the spectrum of a field.
Namely, then $X \to \Spec(K)$ will be a morphism of schemes and
we conclude by the equivalence of radicial and universally injective for
morphisms of schemes, see
Morphisms, Lemma \ref{morphisms-lemma-universally-injective}.

\medskip\noindent
Let us prove (1). Assume $f$ is decent and universally injective. By
Decent Spaces,
Lemmas \ref{decent-spaces-lemma-base-change-relative-conditions},
\ref{decent-spaces-lemma-composition-relative-conditions}, and
\ref{decent-spaces-lemma-properties-trivial-implications}
(to see that an immersion is decent) we see that the discussion in
the first paragraph applies.
Let $X$ be a nonempty decent reduced algebraic space
universally injective over a field $K$. In particular we see that $|X|$
is a singleton. By
Decent Spaces, Lemma \ref{decent-spaces-lemma-when-field}
we conclude that $X \cong \Spec(L)$ for some extension
$K \subset L$ as desired.

\medskip\noindent
A quasi-separated morphism is decent, see
Decent Spaces,
Lemma \ref{decent-spaces-lemma-properties-trivial-implications}.
Hence (1) implies (2).

\medskip\noindent
Let us prove (3).
Recall that the separation axioms are stable under base change
and composition and that closed immersions are separated, see
Morphisms of Spaces,
Lemmas \ref{spaces-morphisms-lemma-base-change-separated},
\ref{spaces-morphisms-lemma-composition-separated}, and
\ref{spaces-morphisms-lemma-immersions-monomorphisms}.
Thus the discussion in the first paragraph of the proof applies.
Let $X$ be a reduced algebraic space universally injective and
locally separated over a field $K$.
In particular $|X|$ is a singleton hence $X$ is quasi-compact, see
Properties of Spaces, Lemma \ref{spaces-properties-lemma-quasi-compact-space}.
We can find a surjective \'etale morphism $U \to X$ with $U$ affine, see
Properties of Spaces,
Lemma \ref{spaces-properties-lemma-quasi-compact-affine-cover}.
Consider the morphism of schemes
$$
j :
U \times_X U
\longrightarrow
U \times_{\Spec(K)} U
$$
As $X \to \Spec(K)$ is universally injective $j$ is surjective,
and as $X \to \Spec(K)$ is locally separated $j$ is an immersion.
A surjective immersion is a closed immersion, see
Schemes, Lemma \ref{schemes-lemma-immersion-when-closed}.
Hence $R = U \times_X U$ is affine as a closed subscheme of an affine scheme.
In particular $R$ is quasi-compact.
It follows that $X = U/R$ is quasi-separated, and the result follows from (2).
\end{proof}

\end{document}
